\section{Approach}\label{sec:approach}\label{sec:system}

\inbrief{A smooth surrogate moves all blocks together, a legalizer removes overlaps,
and coordinate descent keeps a change only if the exact score drops, at milliseconds
per check. Two add-on lanes for escaping stalls, kicks along negative-curvature
directions and optimal reassignment of blocks to slots, were committed on
21--22 May; we infer they were never scored.}

The surrogate proposes and the exact score accepts (\cref{fig:architecture}); this
section describes each stage, then what was scored.

\begin{figure}[tbp]
\centering
\begin{tikzpicture}[font=\small,
  stage/.style={draw=blue!65!black,fill=white,rounded corners=2pt,line width=0.6pt,
    text width=29mm,minimum height=15mm,align=center,inner sep=3pt},
  exact/.style={stage,draw=orange!85!black},
  ext/.style={draw=muted,dashed,fill=white,rounded corners=2pt,line width=0.6pt,
    text width=36mm,minimum height=12mm,align=center,inner sep=2pt},
  arr/.style={-{Latex[length=2mm]},thick,color=muted},
  darr/.style={-{Latex[length=2mm]},thick,dashed,color=muted}]
  \node[stage] (init) {\textbf{Initialize}\\soft rectangles,\\spread apart};
  \node[stage,right=5mm of init] (smooth) {\textbf{Smooth descent}\\Adam on a\\differentiable score};
  \node[stage,right=5mm of smooth] (legal) {\textbf{Legalize}\\largest first,\\zero overlap};
  \node[exact,right=5mm of legal,text width=34mm] (cd) {\textbf{Coordinate descent}\\keep a move only if\\the exact score drops};
  \node[right=3.5mm of cd,align=center,inner sep=1pt] (out) {placement};
  \node[exact,below=6mm of cd,text width=42mm,minimum height=11mm] (eval)
    {\textbf{Incremental evaluator}\\touched nets and cells};
  \draw[arr] (init) -- (smooth);
  \draw[arr] (smooth) -- (legal);
  \draw[arr] (legal) -- (cd);
  \draw[arr] (cd) -- (out);
  \draw[{Latex[length=2mm]}-{Latex[length=2mm]},thick,color=orange!85!black] (cd) -- (eval);
  \node[ext,above=8mm of cd] (lanec) {\textbf{Lane C}: Hungarian\\reassignment, re-polish};
  \node[ext,left=4mm of lanec] (laneb) {\textbf{Lane B}: curvature\\kick, re-polish};
  \coordinate (j) at ($(cd.north)+(0,4mm)$);
  \draw[dashed,thick,muted] (cd.north) -- (j);
  \draw[darr] (j) -| (laneb.south);
  \draw[darr] (j) -- (lanec.south);
  \node[left=3mm of laneb,align=right,font=\footnotesize\itshape,color=muted]
    {added 21--22 May;\\never scored};
  \begin{scope}[on background layer]
    \node[fit=(init)(smooth),fill=blue!7,rounded corners=3pt,inner sep=2mm,
      label={[font=\footnotesize,text=blue!65!black]below:smooth surrogate proposes}] {};
    \node[fit=(cd)(eval),fill=orange!9,rounded corners=3pt,inner sep=2mm,
      label={[font=\footnotesize,text=orange!85!black]below:exact score accepts}] {};
  \end{scope}
\end{tikzpicture}
\caption{Architecture. The ranked entry is inferred to have run the solid pipeline:
the smooth surrogate proposes (blue) and the exact, incrementally computed score
accepts (orange). Dashed lanes B and C were committed before the deadline but, by the
same inference, never scored; the predecessor v1 (\cref{sec:scored}) is not shown.}
\label{fig:architecture}
\end{figure}

\subsection{Make the exact score cheap}\label{sec:cheap}\label{sec:incremental-analysis}

The evaluator caches pin positions, net bounding boxes, cell densities, net routing
contributions and macro allowances. Moving object $i$ changes only its pins, its
nets, and the cells under its footprint and along those nets' routes, so the
evaluator subtracts old contributions, adds new ones, re-reads the tails, and can
revert. On 14 distinct moves from IBM01 and IBM10, repeated in two or three batches,
an update took 1.29--3.67\,ms against 1.40--26.87\,s for full recomputation with the
reference Python evaluator, excluding setup and overlap checks (\measurement):
roughly $10^3$--$2\times10^4$ times faster per sampled move. This is what makes
exact search affordable: a 600\,s pass can try about $10^5$ candidates instead of
22--430 with full recomputation. The ratio reflects both the local-update strategy
and a faster implementation, and re-selecting the top $k$ of $G$ cells still costs
time linear in the grid size. Floating-point updates need not cancel exactly, so the
incremental score $\widehat J$ can drift from $\proxy$; the largest relative
component discrepancy was $2.6\times10^{-15}$.

\begin{proposition}[Acceptance under bounded score error]\label{prop:acceptance}
Let $S_p$ and $S_q$ be the incremental states at candidate endpoints $p$ and $q$,
and suppose $|\widehat J(S_p)-\proxy(p)|\le\epsilon$ and
$|\widehat J(S_q)-\proxy(q)|\le\epsilon$. If $\widehat J(S_q)\le\widehat J(S_p)-\tau$,
then
\begin{equation}
 \proxy(q)-\proxy(p)\le-\tau+2\epsilon.
 \label{eq:acceptance-error}
\end{equation}
In particular, $\tau>2\epsilon$ guarantees a strict decrease of the true score.
\end{proposition}
\noindent Proof in Appendix~\ref{app:proofs}. Coordinate descent uses
$\tau=10^{-9}$, and the largest absolute proxy discrepancy over 36 sampled endpoints
was $\epsilon=4.44\times10^{-16}$, so $\tau/2\epsilon\approx1.1\times10^6$
(\measurement): on those states the error would have to be a million times larger
before an accepted move could raise the true score. This is a sampled margin, not a
uniform bound. The sample exercises the single-move path; coordinate descent scores
candidates through a batched float64 path that was not sampled separately, and
rescoring 34 saved endpoints reproduced their recorded scores only to about
$10^{-7}$.

\subsection{Smooth proposals}\label{sec:smooth}

The smooth stage minimizes a differentiable surrogate of \eqref{eq:canonical},
\begin{equation}
\surrogate(p)=\widetilde W_\gamma(p)+\tfrac12\widetilde D(p)+\tfrac12\widetilde C(p)
 +\lambda_t P_{\mathrm{ov}}(p)+\mu P_{\mathrm{out}}(p),
\label{eq:surrogate}
\end{equation}
which replaces each kink of the score by a smooth approximation and adds penalties
for hard-macro overlap ($P_{\mathrm{ov}}$) and for leaving the canvas
($P_{\mathrm{out}}$). It starts from a soft-rectangle initializer that spreads the
hard macros apart by gradient steps; soft clusters start from the benchmark's input
coordinates, which derive from a RePlAce placement (Appendix~\ref{app:objective}).

\paragraph{Wirelength: a smooth maximum.} HPWL uses $\max$ and $\min$, which have no
gradient where the extreme pin changes. Log-sum-exp replaces the maximum by a smooth
one \citep{naylor2001lse}: for values $x_1,\dots,x_r$ and smoothing length
$\gamma>0$,
\begin{equation}
 \max_j x_j\ \le\ \gamma\log\sum_{j=1}^r e^{x_j/\gamma}\ \le\ \max_j x_j+\gamma\log r,
 \label{eq:lse}
\end{equation}
because the sum contains $e^{\max_j x_j/\gamma}$ and is at most $r$ times it; as
$\gamma\to0$ it converges to the true maximum, and $\min_j x_j=-\max_j(-x_j)$. A
large $\gamma$ is smooth but biased, a small one accurate but stiff, so we anneal
$\gamma$ from $5\times10^{-3}$ to $5\times10^{-5}$ of the canvas scale.

\paragraph{Density: Gaussian smearing.} A rectangle's share of a cell changes at a
rate that jumps when its edge crosses the cell boundary. We instead spread each
object $i$, hard macro or soft cluster, of area $A_i=w_ih_i$, as a separable Gaussian
with spreads $\sigma_{xi}=\max(w_i/2,\Delta_x/2)$ and
$\sigma_{yi}=\max(h_i/2,\Delta_y/2)$, where $\Delta_x\times\Delta_y$ is the cell size.
With $\Phi$ the standard normal distribution function, its mass in cell
$b=[x_b^-,x_b^+]\times[y_b^-,y_b^+]$ is
\begin{equation}
m_{ib}=A_i
\Bigl[\Phi\Bigl(\tfrac{x_b^+-x_i}{\sigma_{xi}}\Bigr)-
      \Phi\Bigl(\tfrac{x_b^--x_i}{\sigma_{xi}}\Bigr)\Bigr]
\Bigl[\Phi\Bigl(\tfrac{y_b^+-y_i}{\sigma_{yi}}\Bigr)-
      \Phi\Bigl(\tfrac{y_b^--y_i}{\sigma_{yi}}\Bigr)\Bigr],
\label{eq:gaussian-density}
\end{equation}
a smooth function of position; $\widetilde D$ applies the top-decile rule to the cell
densities $\tilde\rho_b=\sum_im_{ib}/(\Delta_x\Delta_y)$. Top-$k$ selection and
canvas clamps keep some kinks. This smoother density is the surrogate change the
development records credit (\cref{tab:development}, row~9).

\paragraph{Congestion: a routing trace.} RUDY spreads a net's demand uniformly over
its bounding box \citep{spindler2007rudy}; the evaluator routes along rows and
columns, and our surrogate follows it. Each source--sink pin pair adds a horizontal
stripe along the source's row and a vertical stripe along the sink's column, with
soft (sigmoid) cell membership so the stripes move smoothly with the pins; the macro
allowance, smoothing and top-5\% tail follow the evaluator.

\paragraph{Overlap, optimizer and legalization.} $P_{\mathrm{ov}}$ sums hard-macro
penetration areas, with a weight $\lambda_t$ that rises over the first 70\% of the
steps, so the surrogate first finds an arrangement and then pushes blocks apart;
$P_{\mathrm{out}}$ penalizes positions outside the canvas, with weight $\mu$. Adam, a
standard gradient optimizer \citep{kingma2015adam}, takes 500 steps.
\emph{Legalization} then turns the almost-legal result into one with zero overlap
while moving blocks as little as possible: in order of decreasing area, each hard
macro keeps its position if free or moves to the first free spot, in a fixed scan
order, on the nearest ring that has one; a push-apart pass and, if needed, a retry
with a stronger overlap weight remove the rest.

\subsection{Exact local search}\label{sec:cd}

Coordinate descent \citep{wright2015coordinate} visits every movable object, hard
macros and soft clusters alike, in random order. For each coordinate it takes the
legal interval around the current position, bounded by the nearest blocking macros
(within one of the intervals of \cref{sec:geometry}), and scores at most ten
candidates there in one batched incremental call: the ends and midpoint, the current
position, the nearest grid-line crossings, and targets suggested by the object's
nets. The best is kept if it lowers the score by more than $\tau=10^{-9}$. A pass
stops at its time cap or, after half of it, once five consecutive sweeps each gain
less than 0.001. Every accepted step lowers the incremental score and, by
\cref{prop:acceptance} with the sampled margin, the exact score of the internal
float64 state. The price is reach: a one-coordinate move changes a pair's
arrangement only at a point of contact, so coordinate descent stays near the
arrangement it starts from, and the smooth start matters (\cref{sec:simpler}).

\subsection{Escaping plateaus}\label{sec:escape}\label{sec:curvature}\label{sec:assignment-analysis}

When coordinate descent stalls, further gains need many blocks to move at once. We
tried two such \emph{proposals}; every candidate is legalized, re-polished by
coordinate descent and kept only if the exact score drops.

\paragraph{(a) Negative curvature.} Near a placement $p$, a smooth function $f$
behaves like its second-order Taylor expansion along a unit direction $v$:
\begin{equation}
 f(p+tv)\approx f(p)+t\,\nabla f(p)^\top v+\tfrac12t^2\,v^\top \mathbf H v,
 \qquad \mathbf H=\nabla^2 f(p).
 \label{eq:taylor}
\end{equation}
Where small moves stop helping, the gradient term is small; if the curvature
$v^\top\mathbf Hv$ is negative, the surface bends down along $v$, and a step of
moderate length lowers $f$. The most negative curvature is the smallest eigenvalue
of $\mathbf H$, a matrix with about 5.2 million entries on IBM01. We never build
it: automatic differentiation gives a Hessian-vector product $\mathbf Hv$ at a small
multiple of the cost of a gradient \citep{pearlmutter1994hvp}, and the Lanczos
method \citep{lehoucq1998arpack} turns a few dozen such products into a small
subspace whose extreme eigenvalues approximate those of $\mathbf H$. That finds the
most negative curvature direction without ever building the matrix.

\begin{proposition}[Signed curvature decrease]\label{prop:curvature}
Let $f$ be twice continuously differentiable on an open set $U$ containing $p$, and
let $v$ be a unit vector with $v^\top\nabla^2 f(p)v\le-\eta<0$. Choose
$s\in\{-1,1\}$ with $s\,\nabla f(p)^\top v\le0$. Fix $t>0$, and suppose the segment
$\{p+usv:0\le u\le t\}$ lies in $U$ and the Hessian is $L$-Lipschitz on it
(Euclidean norm for vectors, induced norm for matrices). Then
\begin{equation}
 f(p+tsv)-f(p)\le-\frac{\eta}{2}t^2+\frac{L}{6}t^3.
 \label{eq:curvature-taylor}
\end{equation}
For $L>0$ and $0<t\le3\eta/(2L)$ this is at most $-\eta t^2/4$; for $L=0$ the bound is
$-\eta t^2/2$ throughout the segment.
\end{proposition}
\noindent Proof in Appendix~\ref{app:proofs}. In words: step along a direction of
negative curvature, with the sign chosen so the slope does not hurt, and $f$ falls by
at least $\eta t^2/4$ as long as the step is too short for the curvature to change
much; the gradient need not vanish. The guarantee does not carry over to the kicks we
ran. Legalization moves the endpoint after the step. The function $f$ was not the
exact score or the surrogate of \cref{sec:smooth} but an older surrogate kept for the
kicks, whose density and congestion terms multiply horizontal by vertical overlap
lengths; such products have negative-curvature directions built in, so finding one
says nothing by itself about the true score. And at the diagnosed IBM01 endpoint that
surrogate changes branch within $10^{-7}\,\mu$m, while the steps are
0.3--3\,$\mu$m long, so the Lipschitz hypothesis fails at the step length used.

\paragraph{(b) Assignment.} Some gains need two blocks to trade places, which
coordinate descent cannot do without passing through an overlap. The assignment move
draws $K=50$ macros at random from the 150 with the highest connectivity scores,
takes 100 candidate slots that include the macros' current centers, so staying put is
always possible, and scores each single move against a frozen background,
$\Delta_{ij}=\proxy(p^{i\leftarrow j})-\proxy(p^0)$, where $p^{i\leftarrow j}$ moves
only macro $i$ to slot $j$. The Hungarian method \citep{kuhn1955hungarian} then
solves
\begin{equation}
 \min_z\sum_{i,j}\Delta_{ij}z_{ij}
 \quad\text{subject to}\quad
 \sum_jz_{ij}=1,\quad\sum_iz_{ij}\le1,\quad z_{ij}\in\{0,1\}
 \label{eq:assignment-model}
\end{equation}
exactly, in time cubic in $K$. The moves are applied in order of predicted gain,
skipping any that would overlap, and the batch is kept only if the incremental score
falls by more than $10^{-7}$ (a sampled margin of $1.1\times10^8$ by
\cref{prop:acceptance}; \measurement).

\begin{figure}[tbp]
\centering
\begin{tikzpicture}[x=0.62cm,y=0.62cm,font=\footnotesize,
  mac/.style={draw=blue!70!black,fill=blue!12,line width=0.6pt,minimum size=4.2mm,inner sep=0pt},
  term/.style={draw=ink,fill=ink,regular polygon,regular polygon sides=3,minimum size=2.6mm,inner sep=0pt},
  shared/.style={draw=orange!85!black,line width=1.6pt},
  tether/.style={draw=muted,line width=0.7pt}]
  \foreach \yy in {2.2,0} {
    \draw[rulegray] (0,\yy-0.55) -- (14,\yy-0.55);
    \foreach \t in {0,2,...,14} {\draw[rulegray] (\t,\yy-0.65) -- (\t,\yy-0.45);}
    \foreach \s in {3,11} {\draw[dashed,muted] (\s,\yy-0.55) -- (\s,\yy+0.45);}
  }
  \foreach \t in {0,2,...,14} {\node[muted,font=\scriptsize] at (\t,-0.95) {\t};}
  \node[muted,font=\scriptsize] at (14.5,-0.55) {$x$};
  \node[anchor=east,align=right] at (-0.4,2.2) {start\\$F=20$};
  \draw[shared] (2,2.35) -- (12,2.35);
  \node[mac] at (2,2.2) {1}; \node[mac] at (12,2.2) {2};
  \node[term] at (2,1.55) {}; \node[term] at (12,1.55) {};
  \node[anchor=east,align=right] at (-0.4,0) {after swap\\$F=34$};
  \draw[shared] (3,0.15) -- (11,0.15);
  \draw[tether] (11,0.3) |- (2,0.62) -- (2,-0.5);
  \draw[tether] (3,0.3) |- (12,0.92) -- (12,-0.5);
  \node[mac] at (11,0) {1}; \node[mac] at (3,0) {2};
  \node[term] at (2,-0.65) {}; \node[term] at (12,-0.65) {};
  \node[muted,font=\scriptsize] at (3,3.0) {slot $3$};
  \node[muted,font=\scriptsize] at (11,3.0) {slot $11$};
  \node[anchor=west,align=left] at (14.9,1.1)
    {$\Delta=\begin{pmatrix}-1&-9\\-9&-1\end{pmatrix}$\\[2pt]
     rows: macros 1, 2\\ columns: slots 3, 11\\[2pt]
     predicted: $-18$\\ actual: $+14$};
  \draw[shared] (0.2,-1.75) -- (1.0,-1.75);
  \node[anchor=west] at (1.1,-1.75) {shared net, weight 2};
  \draw[tether] (6.8,-1.75) -- (7.6,-1.75);
  \node[anchor=west] at (7.7,-1.75) {net to fixed terminal};
  \node[term] at (13.9,-1.75) {};
  \node[anchor=west] at (14.2,-1.75) {terminal};
\end{tikzpicture}
\caption{An optimal assignment can make things worse: each of the two crossing moves
saves 9 on its own, so the matrix $\Delta$ predicts $-18$ for the swap, but done
together the moves cross and wirelength rises from 20 to 34.}
\label{fig:assignment}
\end{figure}

The matrix sees each move alone, but moves interact (\cref{fig:assignment}). Two
macros at $x_1=2$ and $x_2=12$ share a weight-2 net, and each has a unit-weight net to
a fixed terminal at its start, so up to normalization their wirelength is
\begin{equation}
 F(x_1,x_2)=2|x_1-x_2|+|x_1-2|+|x_2-12|.
 \label{eq:assignment-example}
\end{equation}
For slots 3 and 11, the Hungarian method picks $(x_1,x_2)=(11,3)$ with predicted
change $-18$, yet $F(2,12)=20$ and $F(11,3)=34$: the real change is $+14$, an
interaction residual of $+32$, with every placement legal for macros narrower than
one unit. Single-move gains ignore shared nets and, in the full score, the coupling
of all moves through the top-$k$ tails. Independent-set matching avoids the
wirelength part of this error by matching only cells that share no nets
\citep{chen2008ntuplace3}; our lane picks connected macros on purpose and relies on
the exact recheck.

\subsection{What was scored and what was added later}\label{sec:scored}

\paragraph{The ranked pipeline.} We infer that the entry verified at 0.9974,
resubmitted on 20 May, ran the solid pipeline of \cref{fig:architecture}, which we
call the \emph{single pipeline}, with a 720\,s budget, up to 600\,s of it for
coordinate descent. We did not record which commit the organizers ran; we infer the
configuration from the listed 12 minutes per design and the resubmission date
(Appendix~\ref{app:records}). Its smooth descent runs in single precision on a CUDA
GPU, an Apple GPU or the CPU, whichever is available, so the trajectory depends on
the device, and we have no record of the device used for the judged run.

\paragraph{Lanes B and C.} On 21--22 May, before the deadline, we committed a
three-lane system: Lane A is the same pipeline with a longer budget (2,700\,s, up to
900\,s of coordinate descent), Lanes B and C add the curvature kicks and assignment
batches, and the lowest exact score with zero overlaps wins. By the same inference it
was never scored, so 0.9974 says nothing about it.

\paragraph{The predecessor, v1.} Our 13 May entry, v1, polished two initial
placements serially with coordinate descent, large-neighborhood search and simulated
annealing \citep{kirkpatrick1983sa}, then ran a curvature cascade; a 17 May run of a
later v1 variant took about 53 minutes per design.
